\section{Residual-scale reduction and ramification certificates}\label{sec:residual-scale}
The radius in Definition~\ref{def:gap} describes a region on which inverse
branches exist. It need not describe the region explored by a minimizing
residual. This distinction permits a stronger reduction theorem.

\begin{theorem}[Reduction at the residual scale]\label{thm:residual-scale}
Keep the geometric hypotheses preceding Definition~\ref{def:gap}, and suppose
\begin{equation}\label{eq:residual-gap}
 \sup_{\Omega_t}\|A_F\|\le K_0,\qquad
 c(t)/r(t)\longrightarrow0,\qquad c(t)M(t)\longrightarrow0.
\end{equation}
All conclusions of Theorem~\ref{thm:reduction} hold, with its relative error
$K\eps(t)$ replaced by $Kc(t)M(t)$. This includes its exact-zero assertion and
its description of the entire leading residual set. The constants are uniform
under uniform metric bounds, the displayed bound on $A_F$, and uniform limits
in \eqref{eq:residual-gap}.
\end{theorem}
\begin{proof}
Lemma~\ref{lem:param-cover} below supplies the inverse charts without any
derivative-gap assumption. Taylor's formula still gives
$|E_b(U,t)|\le M(t)|U|^2/2$. The trial $U=0$ on any branch has norm at most
$Kc(t)$. Consequently every actual global minimizer has $|U|\le K'c(t)$,
by the first residual coordinate and the lower metric bound. Since $c=o(r)$,
such a minimizer is in the smaller inverse tube. The unrestricted affine
minimizer also has $|U_b^*|\le K'|C_b|\le K'c$, by the formula in the proof
of Theorem~\ref{thm:reduction}. On either set of points, writing
$V=C_b+D_bU$, we have
\[
 \big|\|(U,V+E_b)\|_H-\|(U,V)\|_H\big|
 \le K''M(t)|U|^2\le K'''c(t)M(t)\|(U,V)\|_H.
\]
Use actual minimizers for the lower bound and affine minimizers for the upper
bound. The argument makes no division by $C_b$ or by a possibly zero distance.
A zero again forces $U=0$ and then $C_b=0$.

If the branch residues are not eventually zero, bounded slopes make their
affine costs comparable with $|C_b|$. The selected order is therefore
$\rho=\max_b\rho_b$. A bounded normalized residual has $U=O(t^\rho)$.
Since $t^\rho=O(c(t))$, its normalized remainder is
$O(M(t)t^\rho)=o(1)$. The two inclusions identifying
\eqref{eq:intrinsic-sheets} now follow exactly by taking limits and by testing
$U=t^\rho u$, respectively. Every constant used above has the asserted
uniform dependence.
\end{proof}

\begin{corollary}[Continuation across the radius boundary]\label{cor:false-wall}
If $c/r\to0$, $A_F$ is bounded, and $rM=O(1)$, the affine reduction remains
asymptotically exact. In particular, $rM\asymp1$ is not, by itself, a boundary
at which the metric quotient becomes nonlinear.
\end{corollary}
\begin{proof}
$cM=(c/r)(rM)\to0$, so Theorem~\ref{thm:residual-scale} applies.
\end{proof}

We next give finite geometric certificates for the two versions of the gap.
They use the derivative numerators of the original map rather than a
one-variable graph of a previously computed supremum. Write
\begin{equation}\label{eq:polar-ideals}
 J=\det DF,\quad P=DG\,\operatorname{adj}(DF),\quad
 N_2=(J D_xP-P D_xJ)\operatorname{adj}(DF).
\end{equation}
In the last expression the source derivative index is contracted with the
adjugate. Entrywise differentiation gives
\[
 A_F=P/J,\qquad D_xA_F\,DF^{-1}=N_2/J^3.
\]
This identity retains cancellations in the numerator. Replacing it by a
product of separate norm bounds would generally lose the exact criterion.

\begin{theorem}[Divisorial derivative criterion]\label{thm:divisorial}
Suppose $c$ is not identically zero, and shorten the interval so that $c>0$.
Take the compact closure of the positive-time tube graph, recording $a,b,r,c$
as coordinate functions. There is a finite proper real corner presentation of
its germ at $t=0$ on which, on every retained component with nonzero indicated
numerators,
\begin{equation}\label{eq:divisor-data}
 \begin{aligned}
 t&=u_t z^a,& r&=u_r z^s,& c&=u_c z^h,\\
 J&=u_J z^j,& P&=z^p P_0,& N_2&=z^q N_0.
 \end{aligned}
\end{equation}
The scalar units and the norms of $P_0,N_0$ have positive lower and finite
upper bounds; $z_i\ge0$. Every face used below is accessible from positive
time. Put $\beta_i=p_i-j_i$ and $\gamma_i=q_i-3j_i$.
Then \eqref{eq:residual-gap} is equivalent to the following finite tests on
all retained charts:
\begin{equation}\label{eq:divisor-tests}
 \begin{array}{ll}
 \beta_i\ge0 &\text{for all }i,\\
 h_i-s_i>0,\quad h_i+\gamma_i>0 &\text{if }a_i>0,\\
 h_i-s_i\ge0,\quad h_i+\gamma_i\ge0 &\text{if }a_i=0.
 \end{array}
\end{equation}
The original condition $rM\to0$ is obtained by replacing $h_i+\gamma_i$
by $s_i+\gamma_i$. An identically zero numerator imposes no test for its
corresponding derivative. If $c\equiv0$, the minimum is identically zero.

These tests are unchanged in truth value by further proper real refinement
and by feasible source reparametrization transporting the marked fast map.
On a family with a fixed such presentation and uniform unit bounds, strict
inequalities imply uniform power bounds. They persist under perturbations
preserving the displayed orders and those unit bounds.
\end{theorem}
\begin{proof}
All coordinate graphs are semialgebraic. Apply proper real uniformization to
their compact closure, then resolve the product of $t,r,c,J^2,\sum P_{ab}^2$
and $\sum(N_2)_{abc}^2$ on each component on which those factors are not
identically zero. Prune to the closure of positive time after each proper
map. Lemmas~\ref{lem:factor} and \ref{lem:accessible} give
\eqref{eq:divisor-data}, nonvanishing vector units, and accessible corner
boxes covering the germ. A component where a numerator vanishes identically
has zero derivative there and can be treated separately. No component where
$J$ vanishes identically can meet positive time.

For a monomial ratio $f=u z^d$, with $u$ bounded above and away from zero,
boundedness on the retained boxes is equivalent to $d_i\ge0$ for every
coordinate face. Necessity follows by fixing all other divisor coordinates
strictly positive and approaching that face. Moreover, uniform convergence
$f\to0$ as $t\to0$ is equivalent to $d_i>0$ on every face with $a_i>0$
and $d_i\ge0$ on the remaining faces. Necessity follows from the same
transverse arcs. For sufficiency choose
$\delta=\min_{a_i>0}d_i/a_i>0$. After rescaling the boxes to $z_i\le1$,
$|f|\le C t^\delta$. Faces with $a_i=0$ and a negative exponent would also
contradict boundedness at positive time. Applying these facts to $P/J$,
$c/r$, and $cN_2/J^3$ proves exactly \eqref{eq:divisor-tests}.
The supremum of the last ratio is $cM$, so this is a criterion on the whole
tube, not merely along a selected branch.

Every refinement presents the same boundedness and uniform-limit statements;
hence it has the same answer. Fast derivatives are unchanged by a source
coordinate change, by the chain rule, when the tube and $F$ are transported.
Finally the minimum of the finitely many positive ratios and the uniform
unit bounds give uniform powers in a fixed presentation. The perturbation
assertion follows with these same constants.
\end{proof}
The criterion classifies the derivative certificate, not every possible
mechanism that could yield an affine limit. Its refinement invariance does
not assert that arbitrary perturbations admit one simultaneous resolution.
On a fixed resolution type the admissible order vectors form the explicitly
given relatively open rational polyhedral region. Its faces distinguish
loss of residual localization, slope boundedness, and second-derivative
smallness. Corollary~\ref{cor:false-wall} identifies which apparent radius
faces can be crossed while retaining the affine quotient.

\begin{proposition}[A nonlinear residual at the true scale]\label{prop:nonlinear-wall}
Let $\lambda>0$ and, near $(y,x)=(0,0)$, set
\[
 F(y,x)=(y,x^3+y^2x),\quad a(t)=(t,0),\quad
 G(y,x)=yx^2,\quad b(t)=\lambda t^3.
\]
One may take $r(t)=t^2$ on a sufficiently small time interval. The fast fibre
has a unique real point $(t,0)$, the slopes are bounded on the tube,
$c/r\to0$, and $M(t)\asymp t^{-3}$. The complete normalized residual set
at order three is the closed nonlinear set
\begin{equation}\label{eq:cubic-residual}
 E_\lambda=\{(u,\xi^3+\xi,\xi^2-\lambda):u,\xi\in\R\}.
\end{equation}
For every positive metric $H$, the squared leading cost is
$\min_{E_\lambda}y^{\mathsf T}Hy$. In the Euclidean metric it equals
\begin{equation}\label{eq:cubic-cost}
 \min_{z\ge0}\{z^3+3z^2+(1-2\lambda)z+\lambda^2\}.
\end{equation}
For $\lambda>1/2$ this is strictly less than the affine-centre prediction
$\lambda^2$.
\end{proposition}
\begin{proof}
$J=3x^2+y^2$ vanishes only at the origin. The critical value is therefore
zero, and boundary values stay away from the central target after choosing
a compact neighbourhood whose central fibre is interior. The equation
$x^3+y^2x=0$ at $y=t>0$ has only $x=0$. On the tube $y=t+O(t^2)$.
Direct differentiation gives
\[
 A_F=\left(\frac{3x^2(x^2-y^2)}{3x^2+y^2},
                \frac{2xy}{3x^2+y^2}\right).
\]
Its entries are bounded. In coordinates $y=t\eta$, $x=t\xi$, with
$\eta$ bounded above and away from zero, the three distinct second fast
derivatives are
\begin{align*}
 &-\frac{12t\xi^2\eta(9\xi^4+4\xi^2\eta^2-\eta^4)}
              {(3\xi^2+\eta^2)^3},\\
 &\frac{6\xi(\xi^2+\eta^2)(3\xi^2-\eta^2)}
              {t(3\xi^2+\eta^2)^3},\qquad
 -\frac{2\eta(3\xi^2-\eta^2)}{t^3(3\xi^2+\eta^2)^3}.
\end{align*}
The rational factors are bounded for all real $\xi$. At $\xi=0,\eta=1$
the last derivative is $2t^{-3}$. Thus $M\asymp t^{-3}$ and
$cM\asymp1$, whereas $c/r=\lambda t$.

A bounded residual divided by $t^3$ forces $y=t+t^3u$ with bounded $u$.
The second fast coordinate then forces $x=t\xi$ with bounded $\xi$, since
$\xi^3+(y/t)^2\xi$ is proper and strictly increasing. Substitution gives
\eqref{eq:cubic-residual}, and conversely every bounded $(u,\xi)$ is
realized by this substitution. These same coercive bounds localize every
bounded-cost sequence, proving convergence of minima for each $H>0$.
For $H=I$, minimize first over $u$ and put $z=\xi^2$. Expansion gives
\eqref{eq:cubic-cost}; its derivative at zero is $1-2\lambda$.
The minimum is positive, because $\xi^3+\xi=0$ forces $\xi=0$, where the
last coordinate is $-\lambda$. Thus the order really is three.
\end{proof}
The nonlinear set in \eqref{eq:cubic-residual} has the same intrinsic meaning
as the affine sets in \eqref{eq:intrinsic-sheets}: it is the zero-time section
of the normalized feasible residual image. Its metric tensor envelope is
still $\mathcal C(E_\lambda)$. No finite affine or semidefinite description
of that nonlinear envelope is presumed.
